\section{이론}\label{sec:theory}

이 절의 목적은 하나의 공격적 명제를 수학적으로 정리하는 데 있다. 멀티-툴에서 문제는 ``저랭크 steering 자체''가 아니라 ``정지 K-style steering''이다. 본문에 필요한 이론은 세 가지뿐이다. 첫째, 정지 K는 방출된 도구 상태를 조건으로 삼을 수 없으므로 멀티-툴에서 구조적 반복 편향을 갖는다. 둘째, K-측과 Q-측 개입은 한 스텝에서는 같은 연산자족이지만, KV-cache 재사용 때문에 여러 스텝에서는 달라진다. 셋째, signed Q의 부호는 $\beta=0$ 근방에서만 국소적으로 해석할 수 있다. 이 세 결과가 본문의 중심 주장인 ``early-K / global-Q schedule''과 ``Q-only는 regime-specific variant''를 닫는다. 보다 기술적인 perturbation bound는 부록으로 미룬다.

단일 헤드 self-attention을 고려한다(멀티헤드는 헤드 직합). 레이어 $\ell$에서 질의 $q \in \mathbb{R}^d$, key 행렬 $K \in \mathbb{R}^{T\times d}$, value 행렬 $V \in \mathbb{R}^{T\times d_v}$, softmax 가중치 $s_t(q) = \mathrm{softmax}(qK^\top/\sqrt d)_t$, 출력 $o(q) = \sum_t s_t(q) v_t$를 쓴다. \S\ref{sec:method}의 세 연산자는
\begin{equation}\label{eq:ops}
\underbrace{K' = K + \alpha K P_{\mathrm{ont}}}_{\text{정지 K (SEKA형)}},\quad
\underbrace{Q' = Q + \beta Q P_{\mathrm{ont}}}_{\text{signed Q-회전}},\quad
\underbrace{\alpha_\ell = \alpha \mathbf{1}\{\ell < \lfloor L/4\rfloor\},\ \beta_\ell = \beta}_{\text{layer-adaptive K+Q}}
\end{equation}
로 다시 쓰며, 세 연산자가 같은 $B$를 공유하므로 결과 차이는 기저가 아니라 연산자 형식과 레이어 스케줄에서 나온다. 분석 가정은 유계 노름 하나뿐이다: 모든 $q, (k_t, v_t), e_t$에 대해 $\|q\|\le Q_{\max}$, $\|v_t\|\le V_{\max}$, $\|e_t\|\le\rho$. $(q, K, V)$의 결합 분포에 대한 추가 가정은 두지 않는다.

\begin{assumption}[분석 범위]\label{ass:scope}
정리~\ref{thm:qk-duality}와 정리~\ref{thm:beta-star}은 고정된 $(q,K,V)$에서의 per-step 분석이고, 따름정리~\ref{cor:cache-divergence}는 KV-cache가 재사용되는 autoregressive decoding을 따로 다룬다. 따라서 본 절은 ``한 스텝에서 같은가''와 ``여러 스텝에서 같은가''를 명시적으로 구분한다.
\end{assumption}

정지 K는 멀티-툴 상태를 인코딩할 수 없다. 정리~\ref{thm:no-memory}은 이를 ``강한 불가능성''보다 \emph{구조적 불변성}으로 읽는 것이 정확하다. 직관적으로, 같은 hidden state $x$를 만들어내는 두 디코딩 히스토리가 이미 방출한 도구 집합에서 달라도, 정지 연산자는 그 차이를 볼 수 없다.

\begin{theorem}[정지 K는 방출 상태에 무관하다]\label{thm:no-memory}
식~\eqref{eq:ops}의 정지 K-측 연산자를 쓰자. 두 디코딩 히스토리 $h, h'$가 같은 hidden state $x$를 낳지만 방출 도구 집합에서 다르다면 $K'(h) = K'(h')$. 따라서 $(K', K)$의 어떤 함수도 방출 히스토리를 인코딩하지 못한다.
\end{theorem}

\begin{proof}
$K' = K + \alpha K P_{\mathrm{ont}}$는 $(K, \alpha, P_{\mathrm{ont}})$만의 함수이고 셋 다 히스토리 불변이므로 $K' = f(K)$로 쓸 수 있는 히스토리 불변 함수 $f$가 존재한다. $K(h) = K(h')$이면 $f(K(h)) = f(K(h'))$.
\end{proof}

같은 논증이 AdaSEKA에도 적용된다 — 전문가 선택이 현재 입력 $x$의 함수이기 때문이다. 이 구조적 성질의 관측 가능한 귀결은 \texttt{repeated\_first\_tool\_rate}의 상승이며(따름정리~\ref{cor:seka-repeat}, 부록), \S\ref{sec:experiments}이 직접 검증한다. 정리가 말하는 바는 분명하다. 멀티-툴에서 정지 K-only는 원리적으로 방출 상태를 담지 못하므로, 방법론적 개선은 ``더 강한 K''가 아니라 ``K를 어디까지 허용할 것인가''에서 나와야 한다.

Signed Q-회전은 같은 기저 위에서 두 얼굴을 가진다. 정리~\ref{thm:q-sign}이 그 기하를 고정한다: 온톨로지 성분만 $(1+\beta)$배로 조절되고 직교 성분은 보존되므로, $\beta>0$은 온톨로지 정렬 도구의 attention을 증폭하고 $\beta<0$은 감쇠한다.

\begin{theorem}[Q-사영의 기하]\label{thm:q-sign}
식~\eqref{eq:ops}에서 $q' = q + \beta q P_{\mathrm{ont}}$는 $q$의 온톨로지 성분을 $(1+\beta)$배로 균일 축소 또는 증폭하고 직교 성분을 보존한다. $\beta<0$이면 온톨로지 정렬 도구의 attention logit이 감소하고, $\beta>0$이면 증가한다.
\end{theorem}

\begin{proof}
$P_{\mathrm{ont}}^2 = P_{\mathrm{ont}}$에서 $q' P_{\mathrm{ont}} = (1+\beta) q P_{\mathrm{ont}}$, $q'(I-P_{\mathrm{ont}}) = q(I-P_{\mathrm{ont}})$.
\end{proof}

같은 기저 위에서 K-측과 Q-측 개입이 한 스텝 안에서 동등한 작용을 한다는 관찰은 ``같은 연산자족''이라는 서사를 수학적으로 고정한다. 정리~\ref{thm:qk-duality}가 이를 식~\eqref{eq:ops}의 두 축에 대해 직접 진술한다.

\begin{theorem}[Q-K per-step duality]\label{thm:qk-duality}
식~\eqref{eq:ops}의 정지 K-측 연산자 $K' = K + \alpha K P_{\mathrm{ont}}$와 signed Q-회전 $Q' = Q + \beta Q P_{\mathrm{ont}}$를 고려하자. 임의의 한 디코딩 스텝에서 같은 $P_{\mathrm{ont}} = BB^\top$를 공유하고 $\beta = \alpha$이면 두 개입은 attention logit을 동일하게 바꾼다:
\begin{equation}\label{eq:duality}
\frac{1}{\sqrt d}\, q^\top (K')^\top = \frac{1}{\sqrt d}\, (Q' q)^\top K.
\end{equation}
따라서 K-측 강도 $\alpha$와 Q-측 강도 $\beta$는 단일 스텝 attention score 수준에서 서로 교환 가능하다.
\end{theorem}

\begin{proof}
$P_{\mathrm{ont}}^\top = P_{\mathrm{ont}}$이므로 $K' = K(I + \alpha P_{\mathrm{ont}})$의 전치는 $(I + \alpha P_{\mathrm{ont}}) K^\top$이고, 따라서 $q^\top (K')^\top = q^\top (I + \alpha P_{\mathrm{ont}}) K^\top$. 한편 $Q' q = (I + \beta P_{\mathrm{ont}}) q$로부터 $(Q' q)^\top K = q^\top (I + \beta P_{\mathrm{ont}}) K^\top$. 두 식은 $\alpha = \beta$일 때 동일하다.
\end{proof}

정리~\ref{thm:qk-duality}는 두 연산자가 같은 operator family의 두 표상임을 per-step 수준에서 확정하지만 두 표상의 multi-step 거동은 비대칭적이다. 자기회귀 생성에서 key 행렬은 KV-cache에 고정 저장되어 재사용되는 반면 query는 매 스텝 새로 계산되기 때문이다. 따름정리~\ref{cor:cache-divergence}가 이 비대칭을 정식화한다.

\begin{corollary}[KV-cache multi-step divergence]\label{cor:cache-divergence}
KV-cache를 사용하는 autoregressive 디코딩에서 K-측 개입 $K' = K + \alpha K P_{\mathrm{ont}}$은 스텝 $T$에서 캐시에 주입된 이후 모든 후속 스텝 $T+1, T+2, \dots$의 logit에 같은 $\alpha K P_{\mathrm{ont}}$ 항을 누적해 반영한다. 반면 Q-측 개입 $Q' = Q + \beta Q P_{\mathrm{ont}}$는 매 스텝 새 query에만 적용되어 스텝 간 누적이 없다. 따라서 정리~\ref{thm:qk-duality}의 per-step 동등성은 multi-step 해석 수준에서 깨지고, K-측은 long-horizon attention 분포 왜곡을 증폭하며 Q-측은 그러지 않는다.
\end{corollary}

\begin{proof}
캐시된 key는 $k_t^{\text{cache}} = k_t^{\text{orig}} + \alpha P_{\mathrm{ont}}^\top k_t^{\text{orig}}$로 고정되어 이후 모든 스텝의 logit $q_s^\top k_t^{\text{cache}}/\sqrt d$에 $\alpha q_s^\top P_{\mathrm{ont}}^\top k_t^{\text{orig}}/\sqrt d$를 더한다. 따라서 스텝 수 $S$에 대해 K-측 기여는 $\sum_{s=T}^{T+S} \alpha q_s^\top P_{\mathrm{ont}} k_t^{\text{orig}}/\sqrt d$ 형태로 쌓인다. Q-측은 $q'_s = q_s + \beta P_{\mathrm{ont}} q_s$가 각 스텝에서 새로 계산되므로, 스텝 $s$의 개입 항은 $\beta q_s^\top P_{\mathrm{ont}} k_t^{\text{orig}}/\sqrt d$이며 $s$-의존적이지만 스텝 간 누적은 없다.
\end{proof}

정리~\ref{thm:qk-duality}와 따름정리~\ref{cor:cache-divergence}가 묶여 본 논문의 가장 도발적인 해석을 지탱한다. 실패한 것은 ``steering''이 아니라 ``정지 K를 모든 단계에 같은 방식으로 밀어 넣는 배치''다. single step에서는 K-측과 Q-측이 교환 가능하므로 layer-adaptive K+Q와 signed Q-only는 한 스텝 기준에서 같은 가족의 구성원이며, multi-step에서는 K-bias가 캐시에 누적되므로 long-horizon 체제에서 Q-only가 더 안정적인 이유가 설명된다. \S\ref{sec:experiments}은 이 예측을 retail 액션 수 분해와 도메인별 우세 패턴으로 확인한다.

이 결과가 설명해야 할 경험적 사실은 분명하다. retail과 MetaTool에서는 $\beta<0$이 이기고 telecom에서는 $\beta>0$이 이긴다. 어느 부호가 어느 도메인에서 유리한지를 일차 수준에서 예측하는 진단이 다음 정리다. ground-truth 도구의 proxy token 집합을 $\mathcal{G}$라 하고 현재 attention 분포를 $p_0(t) = s_t(q)$라 하자. 온톨로지 부분공간 위의 per-token score를
\begin{equation*}
r_t := \tfrac{1}{\sqrt d}\langle BB^\top q, k_t\rangle
\end{equation*}
로 정의한다. 즉 $r_t$는 $q$의 온톨로지 성분이 각 key와 얼마나 정렬되어 있는지를 재는 값이다. ground-truth에 걸린 attention 질량 $L(\beta) := \sum_{t\in\mathcal{G}} p_\beta(t)$의 $\beta=0$ 근방 변화를 살피면 다음 진단이 나온다.

\begin{theorem}[signed Q의 일차 개선 방향]\label{thm:beta-star}
$\pi_{\mathcal{G}} = \sum_{t\in\mathcal{G}} p_0(t) > 0$, $\bar r = \sum_t p_0(t) r_t$, $\bar r_{\mathcal{G}} = \frac{1}{\pi_{\mathcal{G}}}\sum_{t\in\mathcal{G}} p_0(t) r_t$라 하자. 그러면
\begin{equation}\label{eq:beta-star}
\left.\frac{dL}{d\beta}\right|_{\beta=0} = \pi_{\mathcal{G}}\,(\bar r_{\mathcal{G}} - \bar r).
\end{equation}
특히 $\bar r_{\mathcal{G}} - \bar r > 0$이면 $\beta=0$ 근방에서 $\beta>0$ 방향이, $\bar r_{\mathcal{G}} - \bar r < 0$이면 $\beta<0$ 방향이 ground-truth attention mass를 증가시키는 국소 방향이다.
\end{theorem}

\begin{proof}
softmax gradient identity $\partial_\beta p_\beta(t) = p_\beta(t)(r_t - \mathbb{E}_{p_\beta}[r])$에 $\mathcal{G}$에 대한 합을 취하고 $\beta=0$에서 평가하면 $\frac{dL}{d\beta}\big|_{\beta=0} = \sum_{t\in\mathcal{G}} p_0(t)(r_t - \bar r) = \pi_{\mathcal{G}}(\bar r_{\mathcal{G}} - \bar r)$. 따라서 국소 개선 방향의 부호가 $\bar r_{\mathcal{G}} - \bar r$의 부호와 일치한다.
\end{proof}

이 정리는 signed Q가 어떤 regime에서 focus처럼, 어떤 regime에서 coverage처럼 작동하는지를 한 값 --- ground-truth proxy에 걸린 온톨로지 정렬도의 조건부 평균에서 전체 평균을 뺀 값 --- 의 부호로 환원한다. 하지만 이 결과가 주는 것은 \emph{local direction}이지 전역 최적 $\beta$가 아니다. 따라서 본문은 이 정리를 배포 가능한 router가 아니라 regime 설명 도구로만 사용한다. 네 Qwen 도메인에서의 부호 방향이 실측 우세 방향과 어떻게 대응되는지는 표~\ref{tab:sign-routing}에 정리한다.

\begin{table}[h]
\centering
\caption{정리~\ref{thm:beta-star}의 경험적 진단 (Qwen2.5-7B). $\beta$ 부호에 대한 국소 개선 방향 신호.}
\label{tab:sign-routing}
\small
\begin{tabularx}{\linewidth}{@{}Xcccc@{}}
\toprule
도메인 & baseline F1 & $\operatorname{sign}(\bar r_{\mathcal{G}} - \bar r)$ & 실측 우세 방향 & 해석 \\
\midrule
$\tau^2$ Telecom & 0.251 & $+$ & $\beta=+0.10$, $+22.74$pp & under-focused regime \\
$\tau^2$ Retail & 0.468 & $-$ & $\beta=-0.03$, $+5.11$pp & coverage regime \\
$\tau^2$ Airline & 0.329 & $+$ & layer-adaptive 우세, $\beta<0$은 약함 & mixed boundary \\
MetaTool ST4 & 0.731 & $-$ & $\beta=-0.03$, $+2.28$pp & over-confident regime \\
\bottomrule
\end{tabularx}
\end{table}

정지 K는 왜 모든 레이어에 균일하게 두면 위험한가. 이를 다루는 마지막 결과는 보조적이다. 메인 주장은 앞선 세 결과와 \S\ref{sec:experiments}의 stepwise / action-count 증거만으로 닫히며, 아래 정리는 late-layer K perturbation이 더 취약할 수 있음을 뒷받침하는 정량적 부록 결과로 읽어야 한다.

\begin{theorem}[Attention 출력의 exact integral remainder]\label{thm:bound}
$\phi(\tau) := \sum_t \mathrm{softmax}(\ell(q)+\tau\alpha(q))_t v_t$에 대해 $\hat o(q) - o(q) = L(q, E) + R(q, E)$가 정확한 등식으로 성립하고, 유계 노름 아래
$\|\hat o - o\|^2 \le 2\,\mathrm{Var}_s\alpha(q)\cdot\mathrm{Var}_s[V](q) + C_1 \rho^4$, $C_1 = 2 Q_{\max}^4 V_{\max}^2/d^2$.
\end{theorem}

본문 논증은 한 줄 Taylor 전개와 softmax Jacobian으로 충분하며, 완전 증명은 부록~\ref{app:proof:bound}에 있다. \S\ref{sec:experiments}의 실측은 Qwen L=13에서 좌변/우변 비율의 중앙값이 $2.36\times 10^{-8}$로 bound가 공허하지 않음을 확인한다. 이 보조 결과가 말해 주는 바는 단순하다. 같은 $\alpha$라도 late-layer K perturbation은 더 위험할 수 있고, early-layer K perturbation은 정보 각인 쪽에 더 유리하다. 본 논문의 스케줄 $\alpha_\ell = \alpha \mathbf{1}\{\ell<\lfloor L/4\rfloor\}$, $\beta_\ell = \beta$는 이 설계 원리의 최소 구현이다.

다섯 결과는 같은 연산자족 \eqref{eq:ops}의 서로 다른 축을 짚는다. 정리~\ref{thm:no-memory}은 연산자 정의만으로 history-free 한계를, 정리~\ref{thm:q-sign}은 사영 멱등성만으로 Q 부호의 기하를, 정리~\ref{thm:qk-duality}는 $P_{\mathrm{ont}}$의 대칭성만으로 K와 Q의 per-step 동등성을, 따름정리~\ref{cor:cache-divergence}는 KV-cache 재사용만으로 multi-step 비대칭성을 지적하며, 정리~\ref{thm:beta-star}과 정리~\ref{thm:bound}은 softmax gradient와 Taylor 전개로 각각 국소 방향과 레이어별 교란 경계를 제공한다. 이로써 layer-adaptive K+Q와 signed Q-only는 ``같은 기저, 같은 per-step 작용, 서로 다른 multi-step 거동''이라는 하나의 연산자족 내 두 끝으로 해석된다.
